%% ── 10. EVALUACIÓN DE RESULTADOS ─────────────────────────────────────────────
\chapter{Evaluación de Resultados}

Este capítulo presenta los resultados del análisis del MVP sobre el contrato
piloto Contrato A, la distribución y tipología de hallazgos, la
comparación entre modelos LLM evaluados, los hallazgos de mayor confianza
confirmados de forma independiente, y un análisis crítico de las fortalezas
y limitaciones observadas durante la evaluación.

\section*{Contrato Evaluado: Contrato A}

El contrato seleccionado para la evaluación del MVP es el Contrato A, concesionado por
ProInversión. Sus características lo convierten en un caso de prueba
representativo y exigente:

\begin{table}[H]
\centering
\caption{Características del contrato A}
\label{tab:contrato_ptar}
\begin{tabular}{p{4.5cm} p{8.5cm}}
\toprule
\textbf{Característica} & \textbf{Valor} \\
\midrule
Extensión total           & 324 páginas \\
Número de secciones       & 41 \\
Número de cláusulas       & 448 \\
Total de caracteres       & 749{,}499 \\
Tipo de contrato          & Concesión APP — agua y saneamiento \\
Concedente                & ProInversión (Estado peruano) \\
Estructura                & Tablas anidadas, anexos técnicos, numeración mixta \\
Complejidad normativa     & D.Leg.\ 1362, Reglamento APP, normas SUNASS \\
\bottomrule
\end{tabular}
\end{table}

Este contrato fue elegido por su alta complejidad estructural (tablas
dentro de tablas, secciones con sub-numeración no estándar) y su riqueza
normativa (múltiples referencias cruzadas a cuerpos legales sectoriales),
lo que representa el escenario más exigente para un sistema de auditoría
automática.

\section*{Resultados Cuantitativos por Modelo}

\begin{table}[H]
\centering
\caption{Resultados comparativos del MVP — Contrato A (324 páginas)}
\label{tab:resultados_mvp}
\begin{tabular}{p{4.5cm} p{3cm} p{3cm}}
\toprule
\textbf{Métrica} & \textbf{Gemini 2.5 Pro} & \textbf{Gemini 3.1 Preview} \\
\midrule
Total hallazgos           & \textbf{119}       & 67 \\
Severidad ALTA            & 68 (57\%)          & 38 (57\%) \\
Severidad MEDIA           & 39 (33\%)          & 23 (34\%) \\
Severidad BAJA            & 12 (10\%)          & 6 (9\%) \\
Secciones auditadas       & 41/41 (100\%)      & 41/41 (100\%) \\
Tiempo total              & $\approx 75$ min   & $\approx 170$ min \\
Tiempo promedio/sección   & $\approx 45$ seg   & $\approx 3$--$7$ min \\
Nodos en grafo (GraphRAG) & 208                & 477 \\
Relaciones en grafo       & 258                & 617 \\
\bottomrule
\end{tabular}
\end{table}

\section*{Distribución de Hallazgos por Tipo}

Los hallazgos se clasifican en cuatro categorías según la naturaleza
de la inconsistencia detectada:

\begin{table}[H]
\centering
\caption{Distribución de hallazgos por tipo — Gemini 2.5 Pro}
\label{tab:hallazgos_tipo}
\begin{tabular}{p{4cm} p{2cm} p{2cm} p{5cm}}
\toprule
\textbf{Tipo} & \textbf{Count} & \textbf{\%} & \textbf{Descripción} \\
\midrule
ERROR\_PLAZO &
    34 & 28.6\% &
    Contradicciones entre fechas, plazos o cronogramas en distintas secciones \\
INCONSISTENCIA &
    31 & 26.1\% &
    Incoherencias semánticas internas: obligaciones contradictorias o
    redacciones mutuamente excluyentes \\
ERROR\_LOGICO &
    27 & 22.7\% &
    Fallas lógicas: imposibilidades cronológicas, circularidades, condiciones
    que nunca pueden cumplirse \\
OMISION\_NORMATIVA &
    27 & 22.7\% &
    Ausencia de requisitos exigidos por D.Leg.\ 1362, Reglamento APP o
    normativa sectorial, detectados por el módulo RAG \\
\midrule
\textbf{Total} & \textbf{119} & \textbf{100\%} & \\
\bottomrule
\end{tabular}
\end{table}

La distribución es relativamente uniforme entre los cuatro tipos, lo que
indica que el sistema cubre de manera equilibrada las distintas dimensiones
de análisis (temporal, semántica, lógica y normativa) sin sesgo hacia un
único tipo de inconsistencia.

\section*{Distribución de Hallazgos por Agente}

Cada hallazgo fue generado por uno de los tres agentes especializados:

\begin{table}[H]
\centering
\caption{Hallazgos por agente y área de especialización — Gemini 2.5 Pro}
\label{tab:hallazgos_agente}
\begin{tabular}{p{2.5cm} p{2cm} p{8cm}}
\toprule
\textbf{Agente} & \textbf{Hallazgos} & \textbf{Especialización principal} \\
\midrule
Jurista  & 47 (40\%) & Omisiones normativas; inconsistencias con D.Leg.\ 1362
                       y guías ProInversión \\
Auditor  & 52 (44\%) & Contradicciones internas; referencias cruzadas rotas;
                       incoherencias entre secciones \\
Cronista & 20 (17\%) & Conflictos de fechas; plazos inconsistentes;
                       imposibilidades cronológicas \\
\midrule
\textbf{Total} & \textbf{119} & \\
\bottomrule
\end{tabular}
\end{table}

El Auditor generó la mayor proporción de hallazgos (44\%), lo que refleja
la alta densidad de referencias cruzadas en un contrato de concesión APP de
esta complejidad. El Jurista aportó el 40\% con omisiones e inconsistencias
normativas, confirmando el valor del módulo RAG para contextualizar el
análisis con la legislación vigente.

\section*{Hallazgos de Alta Confianza (Confirmados por Ambos Modelos)}

Los siguientes hallazgos fueron identificados de forma independiente por
Gemini~2.5~Pro y Gemini~3.1~Pro~Preview, lo que les otorga el mayor nivel
de confianza al descartar variabilidad del modelo:

\begin{enumerate}
    \item \textbf{Cláusula 3.3.h} — \texttt{ERROR\_PLAZO} \\
    \textit{Descripción:} Contradicción en los plazos del procedimiento de
    transferencia de Participación Mínima. La cláusula 3.3.h establece un
    plazo de 30 días que es incompatible con el plazo de 15 días requerido
    por el anexo técnico de transferencia. \\
    \textit{Impacto potencial:} Ambigüedad jurídica en procesos de transferencia
    accionaria; riesgo de disputa entre las partes.

    \item \textbf{Cláusula 4.7} — \texttt{ERROR\_PLAZO} \\
    \textit{Descripción:} La suspensión de 180 días establecida en 4.7 contradice
    los plazos de notificación y reactivación de las Cláusulas 4.13 y 4.14,
    que presuponen una suspensión máxima de 90 días. \\
    \textit{Impacto potencial:} Riesgo operativo en situaciones de fuerza mayor;
    el concesionario podría encontrarse en incumplimiento sin justificación
    clara.

    \item \textbf{Cláusula 4.13} — \texttt{ERROR\_LOGICO} \\
    \textit{Descripción:} El mecanismo de aprobación tácita está mal definido:
    la condición de silencio positivo puede activarse antes de que el regulador
    haya recibido formalmente la documentación completa. \\
    \textit{Impacto potencial:} Aprobaciones automáticas de obras o modificaciones
    sin revisión efectiva del concedente.

    \item \textbf{Cláusula 6.36.b} — \texttt{ERROR\_LOGICO} \\
    \textit{Descripción:} La cláusula exige observaciones sobre una etapa
    constructiva que, según el cronograma del contrato (Anexo 3), ya habría
    concluido antes de la fecha de suscripción. \\
    \textit{Impacto potencial:} Obligación contractual que resulta físicamente
    imposible de cumplir; puede ser usada para argumentar incumplimiento
    injustificado.

    \item \textbf{Cláusula 8.2} — \texttt{INCONSISTENCIA} \\
    \textit{Descripción:} La fórmula de retribución al concesionario aplica
    deducciones por indicadores de calidad sobre componentes de inversión
    que, según la cláusula 8.1, están excluidos del mecanismo de penalidades
    operativas. \\
    \textit{Impacto potencial:} Error financiero estructural; podría generar
    subpagos recurrentes o reclamos arbitrales.
\end{enumerate}

\section*{Análisis del Grafo de Conocimiento (GraphRAG)}

El módulo GraphRAG construyó un grafo dirigido del contrato con las
siguientes características:

\begin{table}[H]
\centering
\caption{Estadísticas del grafo de conocimiento por modelo}
\label{tab:grafo}
\begin{tabular}{p{5cm} p{3cm} p{3cm}}
\toprule
\textbf{Estadística} & \textbf{Gemini 2.5 Pro} & \textbf{Gemini 3.1 Preview} \\
\midrule
Nodos totales (sujetos/objetos)  & 208 & 477 \\
Relaciones totales               & 258 & 617 \\
Tipo más frecuente               & \texttt{referencia} & \texttt{referencia} \\
Relaciones \texttt{contradice}   & 12  & 31 \\
Relaciones \texttt{obliga\_a}    & 58  & 134 \\
Relaciones \texttt{modifica}     & 24  & 67 \\
Nodos con $\geq 5$ conexiones    & 18  & 43 \\
\bottomrule
\end{tabular}
\end{table}

Los nodos con mayor grado de conectividad ($\geq 5$ relaciones entrantes o
salientes) corresponden a cláusulas centrales del contrato: la cláusula de
retribución (8.2), el régimen de garantías (Cap.\ 7), y las obligaciones de
construcción y operación (Cap.\ 6). Esto es consistente con el análisis
experto: las secciones financieras y las de obligaciones operativas son
los focos de mayor complejidad contractual.

La relación \texttt{contradice} fue la más significativa para la
generación de hallazgos: el 83\% de los \texttt{ERROR\_LOGICO} detectados
por Gemini~2.5~Pro tuvieron correspondencia con una arista
\texttt{contradice} en el grafo.

\section*{Análisis Crítico: Gemini 2.5 Pro vs.\ Gemini 3.1 Pro Preview}

\begin{table}[H]
\centering
\caption{Análisis comparativo de los modelos evaluados}
\label{tab:comparacion_modelos}
\begin{tabular}{p{3.5cm} p{4.5cm} p{4.5cm}}
\toprule
\textbf{Dimensión} & \textbf{Gemini 2.5 Pro} & \textbf{Gemini 3.1 Preview} \\
\midrule
Volumen de hallazgos &
    119 (mayor sensibilidad) &
    67 (mayor precisión estimada) \\
Grafo de conocimiento &
    208 nodos / 258 relaciones &
    477 nodos / 617 relaciones ($\times 2.3$) \\
Latencia &
    $\approx 75$ min & $\approx 170$ min ($\times 2.3$ más lento) \\
Costo estimado / análisis &
    $\approx \$50$--$\$70$ &
    $\approx \$120$--$\$180$ \\
Calidad de explicaciones &
    Concisa y estructurada &
    Más detallada y contextualizada \\
Falsos positivos estimados &
    $\approx 13\%$ (4/30 en muestra) &
    $\approx 8\%$ (estimación preliminar) \\
\bottomrule
\end{tabular}
\end{table}

\textbf{Interpretación de la paradoja grafo/hallazgos:}
Gemini~3.1 construye un grafo de conocimiento $\times 2.3$ más rico, pero
genera 44\% menos hallazgos. Esta aparente paradoja se explica por el
comportamiento diferente ante la ambigüedad: Gemini~2.5~Pro tiende a reportar
toda inconsistencia potencial (alta sensibilidad), mientras que Gemini~3.1
usa el grafo más denso para filtrar relaciones que, si bien son tensions
textuales, no constituyen inconsistencias definitivas en el contexto del
contrato completo (mayor precisión). La distribución de severidad es casi
idéntica entre ambos modelos ($\approx 57\%$ ALTA, $\approx 33\%$ MEDIA,
$\approx 10\%$ BAJA), lo que confirma coherencia metodológica en la
clasificación de riesgo.

\textbf{Recomendación:} Gemini~2.5~Pro es el modelo óptimo para producción
por su velocidad y cobertura. Gemini~3.1 es candidato para una segunda
fase de validación de hallazgos de severidad ALTA, dentro de la arquitectura
futura de agente validador descrita en el Capítulo 5.

\section*{Comparación Global con el Proceso Manual}

\begin{table}[H]
\centering
\caption{ContractIA MVP vs.\ proceso manual — resumen de evaluación}
\label{tab:eval_baseline}
\begin{tabular}{p{3.8cm} p{3.8cm} p{4.5cm}}
\toprule
\textbf{Dimensión} & \textbf{Proceso Manual} & \textbf{ContractIA MVP} \\
\midrule
Tiempo total &
    $\approx 320$ h &
    $\approx 75$ min ($-99.6\%$) \\
Cobertura &
    $\approx 60\%$ (muestreo) &
    $100\%$ ($+67\%$) \\
Inconsistencias detectadas &
    $\approx 25$ (estimación) &
    119 ($\times 4.8$) \\
Reproducibilidad &
    Baja (depende del especialista) &
    Alta (temperatura 0.1) \\
Trazabilidad &
    Manual / parcial &
    Automática al 100\% \\
Costo / contrato &
    $\approx \$7{,}200$ &
    $\approx \$50$--$\$100$ ($-98.6\%$) \\
\bottomrule
\end{tabular}
\end{table}

\section*{Síntesis de la Evaluación}

Los resultados de la evaluación permiten extraer cuatro conclusiones
principales:

\begin{enumerate}
    \item \textbf{Viabilidad técnica confirmada:} El MVP procesó el
    contrato Contrato A de forma completa, sin errores de
    ejecución, en $\approx 75$ minutos, superando los objetivos SMART
    establecidos en el Capítulo 3.

    \item \textbf{Alta cobertura y sensibilidad:} Los 119 hallazgos
    detectados frente a $\approx 25$ estimados por el proceso manual
    demuestran que el análisis multi-agente es significativamente más
    exhaustivo y sistemático que la revisión selectiva por muestreo.

    \item \textbf{Trazabilidad como diferenciador institucional:}
    El 100\% de los hallazgos incluyen ubicación exacta, evidencia textual,
    referencia normativa y razonamiento explícito. Este nivel de
    documentación no existe en el proceso manual y es el principal
    atributo valorado por los especialistas de ProInversión.

    \item \textbf{El sistema no reemplaza al especialista:}
    El $\approx 13\%$ de falsos positivos confirmados y los $\approx 8$
    falsos negativos estimados demuestran que el MVP actúa correctamente
    como \textit{screening} previo. Los hallazgos de severidad ALTA
    requieren validación experta antes de constituir observaciones
    formales al contrato.
\end{enumerate}